\documentclass[12pt, letterpaper]{article}
\usepackage{graphicx} % Required for inserting images
%\usepackage[T1]{fontenc}
\usepackage{hyperref}
\usepackage{tasks}
\usepackage{amsmath}
\usepackage{enumitem}

\title{Programy użytkowe \\ Lista 12}
\date{}
\usepackage{polski}
\begin{document}
\maketitle
\begin{enumerate}
    \item Utwórz dokument, w którym zdefiniujesz następujące środowiska: Definicję, Twierdzenie, Lemat.
    \begin{enumerate}[label=\alph*)]
        \item Twierdzenia i Lematy mają wykorzystywać wspólny licznik.
        \item Definicje oraz równania mają mieć licznik odnoszący się do sekcji, wewnątrz której są zdefiniowane, tzn. w Sekcji 1 mamy Definicję 1.1, 1.2 i tak dalej. 
    \end{enumerate}
\item Zdefiniuj własną komendę, przyjmującą przynajmniej jeden argument, może to być na przykład jakieś wyrażenie matematyczne, którego zapis w kodzie \LaTeX{a} zostanie uproszczony, jeśli zastąpi się je własną komendą.
\item Utwórz plik z przykładową bibliografią i dołącz go do dokumentu używając biblatex albo bibtex.
\end{enumerate}

\end{document}
\maketitle
